\section{本章训练}

\begin{workedexample}{用留数计算实积分}
计算
\[
    I=\int_{-\infty}^{\infty}\frac{dx}{x^2+1}.
\]
在上半圆上积分函数 $1/(z^2+1)$。被包围的唯一极点是 $z=i$，其留数为
\[
    \operatorname{Res}_{z=i}\frac1{(z-i)(z+i)}=\frac1{2i}.
\]
半圆弧上的贡献趋于零，所以留数定理给出
\[
    I=2\pi i\cdot\frac1{2i}=\pi.
\]
这个方法之所以有效，是因为复延拓把一个实积分转化为了孤立奇点处的局部信息。
\end{workedexample}

\begin{applicationbox}{势流与二维场}
若 $f(z)=\phi(x,y)+i\psi(x,y)$ 全纯，则 Cauchy--Riemann 方程推出 $\phi$ 与 $\psi$ 都是调和函数，并且它们的等值线正交相交。在二维静电学和理想流体中，$\phi$ 可以表示势函数，$\psi$ 可以表示流函数。共形映射能把简单几何区域搬运到更复杂的区域，同时保持局部角度。
\end{applicationbox}

\begin{applicationbox}{频率响应}
有理传递函数通常通过复平面中的极点和零点来研究。极点位置控制增长、衰减、共振和稳定性；围道方法可以恢复反变换与渐近行为。因此，复分析为线性系统与信号处理提供了自然语言。
\end{applicationbox}

\begin{chapterexercises}
    \item 验证 $f(z)=z^2$ 的 Cauchy--Riemann 方程，并指出其实部与虚部。
    \item 计算 $\int_{|z|=2}\frac{e^z}{z-1}\,dz$。
    \item 求 $1/[z(z-1)]$ 在 $0<|z|<1$ 中有效的 Laurent 级数。
    \item 分类 $\sin z/z^3$ 在 $0$ 处的奇点，并计算留数。
    \item 用辐角原理计算 $z^3-2$ 在 $|z|=2$ 内的零点数。
    \item 解释为什么非常数全纯函数不能在内部点取得模的最大值。
\end{chapterexercises}

\begin{selectedhints}
    \item $u=x^2-y^2$ 且 $v=2xy$；检查 $u_x=v_y$ 和 $u_y=-v_x$。
    \item Cauchy 积分公式给出 $2\pi i e$。
    \item 写成 $1/[z(z-1)]=-z^{-1}(1+z+z^2+\cdots)$。
    \item 展开为 $z^{-2}-1/6+\cdots$；奇点为二阶极点，留数为 $0$。
    \item 三个零点的模都为 $2^{1/3}<2$。
    \item 最大模原理来自 Cauchy 积分公式或平均值性质；若等号成立，则函数局部常数，进而在连通区域上整体常数。
\end{selectedhints}
